% !PRINTED 452
\weberpage{452}
%<<CONT
enthalten sind. Nun haben wir aus (9):
\weq{12}{\begin{gathered}
\omega' = p\omega, \quad A\omega'^2 + Bp\omega' + Cp^2 = 0, \\
\omega' = \frac{\omega + c}{p}, \\
Ap^2\omega'^2 + (B - 2Ac)p\omega' + Ac^2 - Bc + C = 0,
\end{gathered}}
deren Diskriminante $D'$ ist. Unter den Größen (11) werden aber nur diejenigen der Gleichung (3) genügen, für die die entsprechende Gleichung (12) primitiv ist.

Die erste der Gleichungen (12) ist nach unserer Voraussetzung stets primitiv; von den übrigen sind nur die nicht primitiv, für die
\weq{13}{Ac^2 - Bc + C \equiv 0 \pmod{p},}
und von den Größen (11) fallen soviele aus, die nicht Lösungen von (3) sind, als die Kongruenz (13) Lösungen hat. Bezeichnen wir für den Augenblick die Zahl dieser Lösungen mit $t$, so ist die Anzahl der Werte von $v$, die zu demselben Wert von $u$ gehören, $(p + 1 - t)$.

Nehmen wir zunächst $p = 2$, so haben wir, da $A$ ungerade ist, folgende Fälle:
\[
\begin{aligned}
B &\equiv 0, & & & c &\equiv C \pmod{2}, & t &= 1, \\
B &\equiv 1, & C &\equiv 0, & c &\equiv 0, 1 \pmod{2}, & t &= 2, \\
B &\equiv 1, & C &\equiv 1 & & \pmod{2}, & t &= 0,
\end{aligned}
\]
(weil $c^2 - c$ immer gerade ist). Dies läßt sich so zusammenfassen:
\[
\begin{aligned}
D &\equiv 0 \pmod{4}, & t &= 1, \\
D &\equiv 1 \pmod{8}, & t &= 2, \\
D &\equiv 5 \pmod{8}, & t &= 0.
\end{aligned}
\]
Mit Benutzung des Zeichens $(D, p)$ (\S\,85) können wir also setzen:
\weq{14}{t = (D, p) + 1.}
Ist $p$ ungerade, so ist die Kongruenz (13) gleichbedeutend mit
\[
(2Ac - B)^2 \equiv D \pmod{p},
\]
und diese hat eine Lösung ($2Ac - B \equiv 0$), wenn $p$ in $D$ aufgeht, wenn also $(D, p) = 0$ ist, zwei Lösungen, wenn $D$ quadratischer Rest von $p$, also $(D, p) = +1$ ist, und keine Lösung, wenn $D$ quadratischer Nichtrest, also $(D, p) = -1$ ist. Also gilt auch hier die Formel (14).

Nehmen wir daher vorläufig an, was wir gleich genauer untersuchen werden, es seien unter den $p + 1$ Größen
%CONT>>
